\section{Giới thiệu}\label{sec:intro}

Giảng viên các lớp lập trình nhập môn quy mô lớn không thể đọc hết mọi bài nộp sai, nhưng vẫn muốn biết cả lớp đang sai ở đâu để giảng lại. Một câu trả lời phổ biến là gom cụm các bài sai rồi tóm tắt từng cụm cho giảng viên. Trong cách phát biểu thường gặp nhất, mỗi bài nộp là một đối tượng được mô tả bằng các cặp thuộc tính--giá trị ghi lại test nào bị trượt (biểu diễn đối tượng--thuộc tính--giá trị, OAV); các đối tượng được gom bằng K-means hoặc một độ đo Hamming; sau đó quy nạp các luật sản xuất dạng NẾU~\emph{triệu chứng} THÌ~\emph{quan niệm sai}. Cách làm này hấp dẫn: kết quả test có sẵn ở mọi hệ thống chấm tự động, biểu diễn dễ giải thích, và luật đọc lên giống lời khuyên giảng dạy.

Câu hỏi liệu các cụm đó có gom bài theo \emph{cơ chế} lỗi hay không hiếm khi được kiểm chứng, vì một lý do đơn giản: các bộ dữ liệu công khai không có nhãn cơ chế, và thuê chuyên gia gán nhãn hàng nghìn chương trình vừa tốn kém vừa khó tái lập. Khi thiếu nhãn, chất lượng cụm thường được đánh giá bằng chỉ số nội tại như silhouette, hoặc bằng độ khớp của luật thay thế với cụm. Cả hai đều không cho biết hai chương trình trượt cùng test có sai vì cùng một lý do hay không. Hai chương trình có thể in cùng một số sai: một do vòng lặp dừng sớm một phần tử, một do khởi tạo sai biến tích lũy. Ngược lại, cùng một cách hiểu sai có thể làm trượt các test khác nhau tùy phần còn lại của chương trình được viết thế nào.

Bài báo này xây dựng dữ liệu đánh giá còn thiếu mà không cần gán nhãn mới, rồi dùng nó để kiểm định quy trình kinh điển. Điểm xuất phát là: lần nộp sau của chính sinh viên thường sửa được lần nộp sai. Chúng tôi chạy lại {{n_submissions}} bài C90 của một bộ dữ liệu công khai \citep{orvalho2024cpack} trong môi trường cách ly, ghép mỗi bài sai với lần nộp đạt kế tiếp của cùng sinh viên, và dùng delta debugging để thu gọn khác biệt giữa hai bài thành một bản sửa tối thiểu đã kiểm chứng: tập dòng thay đổi nhỏ nhất mà khi áp vào bài sai thì bài đạt mọi test. Một bộ phân loại tất định gán bản sửa này vào một trong mười hai loại mức chương trình, từ biên vòng lặp đến định dạng in. Như vậy nhãn dựa trên bằng chứng đã được thực thi chứ không dựa trên phán đoán, và mọi nhãn đều có thể suy lại. Vì bản sửa thật mất cân bằng và đôi khi trộn nhiều thay đổi, chúng tôi bổ sung một benchmark thứ hai gồm các mutant một thay đổi của bài đã đạt, với loại cơ chế biết trước do cách tạo.

Với hai benchmark này, chúng tôi trả lời bốn câu hỏi. Chữ ký test làm mất bao nhiêu thông tin về cơ chế? Các mô tả không cần nhãn về \emph{cách} output lệch khỏi đáp án có khôi phục cơ chế tốt hơn ở cùng ngân sách phân cụm không? Luật NẾU--THÌ quy nạp bằng Thuật toán Học Quy nạp (ILA) và phiên bản chịu nhiễu ILA-2 \citep{tolun1998ila,tolun1999ila2} có tổng quát sang sinh viên mới và bài tập mới không? Và embedding mã đóng băng gom theo cơ chế hay theo tác giả và chương trình? Mọi giả thuyết, biểu diễn, trọng số và quy tắc thống kê được đăng ký trước khi đánh giá; mọi thay đổi đều được ghi thành sửa đổi có ghi ngày.

Đóng góp của bài báo gồm:
\begin{enumerate}
\item chạy lại C-Pack-IPAs có thể tái lập trong môi trường cách ly, kèm audit oracle so với hệ thống chấm lịch sử (khớp {{audit_cells_pct}}\% trên {{audit_cells}} verdict test) và lần chạy lại độc lập thứ hai (khớp {{repro_cells_pct}}\%);
\item một benchmark dựa trên bản sửa gồm {{n_pairs}} bản sửa tối thiểu đã kiểm chứng của sinh viên ({{n_real_gold}} nhãn đơn cơ chế) và một benchmark kiểm soát gồm {{n_injected}} mutant một thay đổi, dùng chung một taxonomy và một bộ phân loại;
\item OAV độ lệch thực thi, một biểu diễn không cần nhãn và không phụ thuộc bài của các sai lệch output, cùng protocol đánh giá có trần khả phân biệt và phép so với phân hoạch ngẫu nhiên cùng độ mịn;
\item bằng chứng đăng ký trước rằng chữ ký test trộn lẫn cơ chế (từ chúng chỉ phân biệt được tối đa {{real_H1a_pct}}\% bài thật và {{injected_H1a_pct}}\% bài tiêm lỗi), rằng bằng chứng độ lệch output khôi phục cơ chế tốt hơn trên dữ liệu kiểm soát nhưng chưa đáng tin cậy trên bản sửa thật, rằng luật ILA-2 trên các thuộc tính này chuyển được sang bài tập mới, và rằng embedding mã đóng băng gom chương trình theo nguồn gốc.
\end{enumerate}
Chúng tôi công bố runner replay, các benchmark, protocol và mã phân tích. Trong toàn bài, một nhãn mô tả điều đã thay đổi trong chương trình; nó không khẳng định sinh viên đã tin điều gì.
